% SnipTeX ukázkový dokument
% Autor: Rostislav Brož
% Semestrální práce, 2026

% SnipTeX potřebuje Python backend pro to, aby fungoval.
% Před kompilací dokumentu je potřeba ho nainstalovat:

\iffalse
pip install sniptex
\fi

% Pro kompilaci bude potřeba spustit nasledující sérii příkazů,
% bez nich se citace v dokumentu nejspíš nebudou ukazovat správně.

\iffalse
lualatex --shell-escape docs/main.tex
bibtex main
lualatex --shell-escape docs/main.tex
lualatex --shell-escape docs/main.tex
\fi

\documentclass{article}

\usepackage[a4paper,margin=2cm]{geometry}
\usepackage[czech]{babel}
\usepackage{natbib}
\usepackage[hidelinks]{hyperref}

\usepackage{sniptex}

\begin{document}

\begin{titlepage}
    \centering

    {\Large České vysoké učení technické v Praze\par}
    {\Large Fakulta informačních technologií\par}

    \vspace{3cm}

    {\Huge\bfseries Snip\TeX\par}

    \vspace{0.5cm}

    {\Large Semestrální práce - Typografie a \TeX\par}

    \vfill

    \begin{tabular}{rl}
        Autor: & Rostislav Brož \\
        Rok: & 2026 \\
    \end{tabular}

    \vspace{2cm}
\end{titlepage}

\section*{Snip\TeX}

Snip\TeX{} je La\TeX{}ový balíček pro vkládání označených úryvků kódu z lokálních souborů 
nebo vzdálených repozitářů dostupných přes webové odkazy. 
Extrahuje označené části kódu a sází je jednotným stylem, 
čímž podporuje reprodukovatelnou a udržitelnou dokumentaci.

\subsection*{Instalace}

Snip\TeX{} pro parsování zdrojových souborů používá stejnojmenný Python balíček,
který je publikovaný v softwarovém repozitáři PyPI \citep{pypi}. 
Aby Snip\TeX{} fungoval, je potřeba backend nainstalovat následujícím příkazem.

\IncludeCode[firstnumber=1]{docs/example2.py}{comm0}{cpp}{}

\subsection*{Použití}

Snip\TeX{} definuje příkaz \textbackslash{}IncludeCode, který se používá následovně:

\IncludeCode[firstnumber=1]{docs/example2.py}{comm1}{cpp}{}

Nepovinný argument \textit{options} umožňuje uživateli přepsat nastavení formátování kódu. 
V příkazu níže byly použity přepínače \textit{firstnumber}, \textit{fontsize} a \textit{style},
které nastaví číslování řádků od čísla 1, velikost fontu na \nobreak\textit{\textbackslash{}scriptsize}
a styl zvýrazňování kódu na \textit{monokai}. 
Přepínače musí existovat v balíčku Minted \citep{minted} a být oddělené čárkou. Je také důležité zmínit, že bez úpravy
přepínače \textit{firstnumber} Snip\TeX{} čísluje řádky tak, jak jsou číslované ve~zdrojovém souboru.

\IncludeCode[firstnumber=1]{docs/example2.py}{comm2}{cpp}{}

Povinný argument \textit{source} definuje zdroj, odkud se má
úrvek kódu převzít. Argument může být ve formě cesty k lokálnímu souboru, nebo odkazu na webovou stránku.
Snip\TeX{} dokáže tyto možnosti sám rozlišit. Povinný argument \textit{tag} specifikuje 
unikátní klíčové slovo, které označuje jeden úryvek kódu. Tag musí být oddělený přesně jednou mezerou a 
následovat za frázemi \textit{sniptex-start} a \textit{sniptex-end}, 
které úryvek vymezují (viz Úryvek~1).

\IncludeCode{docs/example.py}{demo}{Python}{Úryvek 1: Ukázka používání tagů.}

Povinným argumentem \textit{language} uživatel zvolí programovací jazyk, v jehož stylu má 
být kód v úryvku zvýrazněn. Posledním argumentem je popisek, který může být prázdný. 
Nicméně složené závorky je potřeba psát i pro prázdný popisek.

\subsubsection*{Lokální úryvek}

Pro ukázku toho, jak příkaz \textbackslash{}IncludeCode funguje pro lokální soubory využijeme
soubor \textit{example.py} s Úryvkem~1. V něm je funkce \textit{main()} obalená tagem \textit{tag1}. 
Příkaz pro zobrazení kódu funkce v dokumentu by tedy vypadal takto:

\IncludeCode[firstnumber=1]{docs/example2.py}{comm3}{cpp}{}

Tento příkaz se přeloží na Úryvek 2:

\IncludeCode{docs/example.py}{tag1}{Python}{Úryvek 2: Úryvek vymezený tagem \textit{tag1} z Úryvku 1.}

\newpage

\subsubsection*{Úryvek z webové stránky}

Součástí domovského repozitáře Snip\TeX{}u na GitHubu je i soubor \textit{example2.py}, 
ve kterém se nachází funkce pro extrakci úryvku. Část této funkce je ohraničená tagem \textit{tag2}, 
tu zobrazíme následujícím příkazem:

\IncludeCode[firstnumber=1]{docs/example2.py}{comm4}{cpp}{}

Výsledkem tohoto příkazu je Úryvek 3:

\IncludeCode[]{https://raw.githubusercontent.com/brozrost/sniptex/main/docs/example2.py}{tag2}{Python}{Úryvek 3: Úryvek kódu z webové stránky.}

\subsection*{Architektura}

Snip\TeX{} se skládá ze dvou hlavních částí,
La\TeX{}ového frontend balíčku a externí backend utility napsané v~Pythonu.
Frontend definuje příkaz \textbackslash{}IncludeCode a zajišťuje sazbu výsledných úryvků pomocí balíčku Minted.
Backend je zodpovědný za přístup ke zdrojovým souborům, vyhledávání označených úseků a jejich extrakci do dočasných souborů.
Při překladu dokumentu frontend pomocí shell escape spustí backend utilitu,
předá jí cestu ke zdrojovému souboru a název tagu,
a následně vloží vygenerovaný dočasný soubor do dokumentu.

\subsection*{Backend}

Python balíček \textit{sniptex} je jednoduché command line rozhraní,
které přijímá zdrojový soubor, název tagu a~cestu k výstupnímu souboru. 
Rozhraní načte zdrojový soubor, vyhledá odpovídající dvojici značek
\textit{sniptex-start} a~\textit{sniptex-end} a následně extrahuje text mezi nimi.

Balíček také provádí kontrolu chybových stavů. Chybovým stavem je myšlen například
neexistující zdrojový soubor, chybějící tag nebo vícenásobný výskyt stejného tagu ve zdrojovém souboru.
V případě detekce chyby backend vyvolá výjimku, která přeruší překlad dokumentu
a vypíše odpovídající chybovou zprávu.
Pro práci s lokálními soubory balíček používá modul \textit{pathlib},
pro vyhledávání tagů modul \textit{re}
a pro načítání vzdálených souborů přes HTTP knihovnu \textit{requests}.

\subsection*{Omezení}

\begin{itemize}
    \item Snip\TeX{} vyžaduje překladač LuaLaTeX.
    
    \item Při kompilaci dokumentu musí být povolen shell escape
    pomocí přepínače \textit{--shell-escape}.

    \item Python backend musí být nainstalovaný a dostupný v systémové proměnné PATH
    pod příkazem \textit{sniptex}.

    \item Balíček závisí na balíčku Minted.

    \item Vzdálené zdrojové soubory musí být dostupné přes HTTP(S)
    a ve formátu čitelného textového souboru.
\end{itemize}

\bibliographystyle{unsrt}
\bibliography{docs/citace.bib}

\end{document}